\documentclass[11pt]{article}
\usepackage{amsmath}
\usepackage{amssymb}
\usepackage{amsthm}
\usepackage{hyperref}
\usepackage{booktabs}

\newtheorem{lemma}{Lemma}[section]

\title{Corrected Policy-Gradient Update Rules for Reinforcement Learning}
\author{Jonathan Gu \\ \texttt{jonathangu@gmail.com} \\ UCLA Department of Economics}
\date{July 2016}

\begin{document}
\maketitle

\begin{abstract}
This article proves that the current consensus on the update rule used in reinforcement
learning does not actually maximize the correct objective. Here we provide the corrected
update rule that trains the parameters of a policy function such that the value of every state
is maximized.
\end{abstract}

\section{Motivation}
Neural networks are used to solve tasks that are hard to solve by ordinary rule based program-
ming. They are used when analytic functions either do not exist or are infeasible to express.
Neural nets are often used to approximate policy functions to determine the best action to take
given a current state, or they are used to approximate value functions to determine the expected
value of being in a given state.
In most economic models it is common specify either a utility of a welfare function. The goal
of agents within the model is to maximize their corresponding function - here we will rename
this function to ``reward function.''
In applied microeconomics researcher commonly use dynamic programming to solve for value
functions by back solving from the final time period. In macroeconomics the common practice
is to use shooting algorithms. In microeconomic theory researchers tend to specify theoretical
models and solve explicitly for some feature of their game. All of the above models specify a
reward function, and therefore it is applicable to frame these problems such that neural nets
become a viable method for approximating policy functions and value functions.
Given the reward function, the state space, and the possible actions, the next step is to solve
for an equilibrium optimal policy function --- which tells what the optimum response to a given
state is --- assuming that all other players are also responding optimally. Here we will refer to
equilibrium play and optimal policy as ``perfect play.''
In the application of games, the economist assumes perfect play. Perfect play is the traditional
notion of equilibrium where everyone is making the best move possible at every time. In games
of immense computational complexity --- such as Go or chess --- perfect play (or equilibrium play)
may not be feasible for the common layman or even for the professional competitor. There are
too many possible future states for a human to process all at once.
The extent of the bounded rationality of the human mind is difficult to measure, but there are
many games where it may be possible to use reinforcement learning to approximate the optimum
value function and the optimum equilibrium moves even when no one knows what such optima
are supposed to look like.
When one trains a neural net --- one starts with supervised learning. The neural net observes
a dataset labelled with good moves and bad moves. The neural net can accordingly develop a
policy function to maximize winning (which will be mathematically formalized below.)
With complex games, and the number of possible future states becomes intractable, the
designer lacks labelled good move/ bad move data sets. In order to approximate the unknown
optimum, one must turn to reinforcement learning, in which the computer simulates a sequence
of moves and discovers a policy function to maximize winning on its own. Reinforcement learning
has already been implemented correctly for a computer to handily defeat world Go champion
Lee Sedol four games to one in March of 2016.
The failure of the human mind to maximize anything has never stopped economists before.
Here we will analyze the how to use reinforcement learning to find optimum equilibrium behavior.
This approach will come in most useful when optimum equilibrium behavior is unknown.

\section{Goal of Article}
In this article, we will show a lemma of reinforcement learning for finite time games.
The current state of the art in neural networks uses supervised learning to train a beginner
policy function, then uses reinforcement learning to further improve the parameters of its approximation.
As discussed in the previous section, supervised learning relies on a labelled dataset to
update its parameters. We will not go into further detail on supervised learning in this article.
Reinforcement learning requires a reward function to tell the value of each state. In modern implementation
of reinforcement learning, parameters are updated according to Ronald J.\ Williams (1992) \cite{williams1992}:
\begin{equation}\label{eq:williams}
\Delta \rho \propto \frac{\partial \log P_\rho(a_t|s_t)}{\partial \rho} z_t
\tag{1}
\end{equation}
Where the goal is to discover parameters corresponding to an optimum value function which
evaluates the value of each state under the assumption of perfect play. The notation will be
described in more detail later, but here we define $\rho$ to be the parameters, $P_\rho$ as the policy
function, $a$ is an action, $s$ is a state, and $z$ is the final outcome. (Conventionally $z$ is positive for
a ``good'' outcome, and negative for a ``bad'' outcome.)
The above rule has all the right ingredients: we update the parameters $\rho$ according to how
the parameters affect the policy function, and if the result is a win then we increase the likelihood
of a similar move later, if the result is a loss then we decrease the likelihood of a similar move.
Although this rule seems to make sense we will show in a lemma below that is is more accurate
to use the slightly different update rule:
\begin{equation}\label{eq:correct}
\frac{\partial v_\rho(s_t)}{\partial \rho}
=
E\!\left[
z_T\sum_{l=t}^T \frac{\partial\log P_\rho(a_l|s_l)}{\partial \rho}
\right]
\tag{2}
\end{equation}
where the expectation is taken according to the policy function $P_\rho$.
The key difference is the expectation which appears in \eqref{eq:correct}. When updating the parameters
of the proposed optimum value function we should look ahead towards the end of the game in a
manner ratifying the proposed optimum value function as correct. \eqref{eq:williams} simply multiplies the final
gradient by $z$ for seemingly no reason.
Furthermore, here we clearly define what we are maximizing: we want to choose the parameters
$\rho$ (which decide the actions) to maximize the value of each state. In other words we want
our policy function to tell us to do actions which will give us the highest reward.

\section{Suitability of Neural Nets}
If we are trying to approximate the continuous function $f:\mathbb{R}^n\to\mathbb{R}$, we can define a single layer
feedforward network $g$ as follows:
\[
g(x)=\sum_{j=1}^n \theta_j h(x^\prime w_j + a_j),\quad a_j,\theta_j\in\mathbb{R},\quad w_j\in\mathbb{R}^n,\ w_j\neq 0,\ n\in\{1,2,3,\dots\}
\]
Where the linear function $x^\prime w_j + a_j$ is the first layer (it is most often a linear function), the
functions $h$ is the single ``hidden'' layer. We can define $h$ to be any function that is $L_p$ integrable
for all $1\leq p\leq\infty$ and squashes: $\lim_{x\to\infty}h(x)=1$ and
\[
\lim_{x\to-\infty}h(x)=0.
\]
A theorem by Hornik, Stinchcombe, and White (1989) \cite{hornik1989} shows that for any triple:
$(\epsilon>0,\ \text{probability measure }\mu,\ \text{compact set }K\subset\mathbb{R}^n)$ there exists $\theta_j,a_j,w_j,n$ such that
\[
\sup_{x\in K}\left|f(x)-g(x)\right|\le \epsilon
\quad\text{and}\quad
\int_K |f(x)-g(x)|\,d\mu\le \epsilon.
\]
Thus even a single layer neural net can approximate any continuous function arbitrarily well.

\section{Overview of Neural Networks}
We analyze games that 1) end in finite time and 2) only reward at the end of the game. We
have the value function $v^*(s)$ which determines the expected value of the outcome of the game
from any state $s$, assuming perfect play by all players.
Examples of real games that have been successfully solved to world championship caliber
with neural nets are: chess (Murray Campbell, Joseph Hoane Jr., Feng-hsiung Hsu (2002) \cite{campbell2002}), checkers
(Jonathan Schaefer et al.\ (1992) \cite{schaeffer1992}), othello (Michael Buro (1999) \cite{buro1999}), and
Go \cite{silver2016}.
In general players do not need to alternate moves, there could be moving simultaneously.
It is not even necessary for there to be two players, there can be any number players in general.
(There could even be one player as in many macroeconomic optimum planning implementations.)

\subsection{Search Approach}
The brute force way to solve a game is to search through all the moves into the future. Since it
takes too much computation to do an exhaustive search through all possible future states, there
are two common methods for finding the optimum policy function:
\begin{enumerate}
\item We can simply truncate the search at a certain depth $\hat t$ into the future and approximate
the value of such a position: $\tilde v(\hat t)$. Truncation is implemented in chess \cite{campbell2002}, checkers \cite{schaeffer1992}, and
othello \cite{buro1999}.
\item After estimating a simple policy function $P_\rho(a|s)$ which tells which action to take depending
on the state, we can use Monte Carlo rollouts to search to the end of the game without
branching into different possible actions. Monte Carlo rollouts are important for this article
because we have a term $E[z_t\mid s=s_t]$ which denotes the final result after starting from
state $s_t$. We can estimate $E[z_t\mid s=s_t]$ by summing over the action sequences according
to the policy $P_\rho(a|s)$. Monte Carlo rollouts are implemented in backgammon \cite{tesauro1996}, and in
Scrabble \cite{sheppard2002}.
\end{enumerate}

\section{Reinforcement Learning}
Reinforcement learning occurs when a software optimizes its parameters on its own. The software
is presented with a reward function, and the software then seeks to find a policy function so that
it can maximize its cumulative reward.
Reinforcement learning differs from supervised learning because the software is never pre-
sented with correct inputs or outputs, and sub-optimal actions are never corrected by an outside
source. The software learns on its own. Software is often trained with supervised learning first,
then trained with reinforcement learning afterwards.
A reinforcement learning model looks exactly like a common microeconomic game --- which
consists of players -- types -- actions -- payoffs -- probabilities. Except that types and actions can
be grouped into a notion of the state, and we use observability rules to guide what is observable
from each state.

The basic reinforcement learning model consists of:
\begin{itemize}
\item States: $S=\{s_0,s_1,s_2,\dots\}$.
\item Possible actions for each state: $A(s)\subset\{a_0,a_1,a_2,\dots\}$.
\item A reward function $r:S\to\mathbb{R}$.
\item Transition rules between states: $f:S\times A\to S$.
\item Observability rules: $h:S\to\mathcal O$, where $\mathcal O$ defines the set of observable attributes of the
states $S$.
\end{itemize}
Here we will add the assumption that the game has a definite end:
\medskip
\noindent\textbf{Assumption 1:} There exists $\bar E\subseteq S$ such that the game ends if it reaches
a state $s\in\bar E$. And for every game where $T$ denotes total number of moves to reach the
end of the game, there exists $\bar T$ such that $T<\bar T$. So that all games will definitely end before
$\bar T$ moves.
\medskip
\noindent\textbf{Assumption 2:} $r(s)=0,\ \forall s\notin \bar E$.

\subsection{Further Notation}
We can define a policy function with regards to the parameters $\rho$ as: $P_\rho(a_t|s_t):S\to\Delta(A)$.
It yields a distribution over the actions $A$, conditional on the state: $s_t$.
$\rho$ are the parameters that determine the policy function. It is the objective of the software
to choose $\rho$ such that in the course of the game, the expected reward is maximized.
We can define the optimum value function:
\[
v_\rho(s)=E\left[z_t\mid s_t=s,\ a_t,\dots,T\sim P_\rho\right].
\]
The notation $a_t,\dots,T\sim P_\rho$ formalizes the assumption that $P_\rho$ is an approximation of perfect
play, and that all future actions will be played according to $P_\rho$. We can simulate draws from
$a_t,\dots,T\sim P_\rho$ with Monte Carlo rollouts --- by simulating to the end of the game without branching
into other possible decision paths.
To simplify notation, we will also use $z_t$ corresponding to the final reward of a game that is
currently in state $s_t$. (In zero sum win/loss games of chess or Go, $z_t\in\{1,-1\}$).

\subsection{Sample Problem Settings}
Here I will provide two examples of game settings that are applicable to reinforcement learning:
\subsubsection{Signalling}
A simple Spence signalling game takes two turns \cite{spence1973}. The worker sends a message after observing his
quality. The firm receives the message and then offers a wage.
\begin{itemize}
\item The State Space $S:(q,m,w)\in S$ is the quality of the worker, the signal of the worker, and
a wage offer. We note that the state is a triplet with possibly empty entries.
\item The quality of the worker is
\[
q=
\begin{cases}
H\in\mathbb{R}_+, & \text{with probability }p \\
L\in\mathbb{R}_+, \ L<H, & \text{with probability }1-p
\end{cases}
\]
\item The observability rule: The firm does not see the quality of the worker.
\item The action space: $A(s):a\in A(s)$
\begin{enumerate}
  \item Worker: any nonnegative message $m\in\mathbb{R}_+$.
  \item Firm: any nonnegative wage $w\in\mathbb{R}_+$.
\end{enumerate}
\item The state transition: $s_{t+1}=f(s_t,a_t)$.
Once a worker makes a signal $m$, the signal $m$ is filled into the state triplet. Once a firm
makes an offer $w$, the offer $w$ is filled into the triplet.
\item The reward function at the end of the game is:
\begin{align*}
r_{\text{worker}}(q,m,w)&=w-m,\\
r_{\text{firm}}(q,m,w)&=q-w.
\end{align*}
Both workers and firms do not receive any reward if the game is not over.
\end{itemize}
\subsubsection{Go}
The game of Go is defined by a $19\times19$ grid. Each point on the grid is either filled with a black
stone, a white stone, or is empty. There is a player that places white stones, and a player that
places black stones. The game is over when the board is filled.
\begin{itemize}
\item State Space $S:s\in S$ is any board position. (Any combination of black stones, white
stones, or empty spaces on the board). There are less than $3^{19}$ possible board positions.
\item Action space: $A(s):a\in A(s)$ is any one move possible from the state $s$. One player
places the stone of his color.
\item State transition: $s_{t+1}=f(s_t,a_t)$ determines the new state at time $t+1$ after move $a_t$
at state $s_t$. Other than alternating moves, there is only one other rule of Go: if a group
of white stones is surrounded by black stones without any adjacent blank spots, then the
white stones are removed from the board. (And vice versa for black stones).
\item The player with the most stones on the board at the end of the game wins. The reward
function is
\[
r_i(s)=
\begin{cases}
0, & \text{if game is not over at state }s,\\
1, & \text{if player }i\text{ won at state }s,\\
-1,& \text{if player }i\text{ lost at state }s.
\end{cases}
\]
The reward function formalizes the assumption that both players only care about the final
win/loss outcome. This follows with Assumption 2 above.
\item Accordingly with Assumption 1, the game ends when the entire board is filled. No Go
game in history has ever exceeded 500 moves.
\end{itemize}

\section{The Optimum Parameters}
We note that the goal of reinforcement learning is to find the optimum parameters $\rho^*$ such that
given any state $s$, $P^*_\rho(s)$ yields the distribution over actions which maximizes the expected final
reward. Finding the optimum $\rho^*$ amounts to finding the optimum policy $P^*_\rho$ and finding the
optimum value function $v^*_\rho$.
In our proof, we will derive an analytic expression for the derivative of the optimum value
function $v^*_\rho$, under the assumption that the current $\rho$ is equal to the optimum $\rho^*$, and that the
policy function and the value function are already the optimum values.
In effect we have found a way to verify that the optimum parameters $\rho^*$ are at a local extrema
by checking if the derivative is zero. However we do not have a convexity argument to claim that
parameters $\rho$ close to the optimum $\rho^*$ will yield a derivative which brings the parameters even
closer to the optimum.
This is why it is common to use supervised learning. First we would use real labelled examples
of good moves and bad moves to train the parameters $\rho$ such that they are slightly closer to the
optimum $\rho^*$. If we wish to use reinforcement learning afterwards, we must hope that our current
parameters $\rho$ are close enough to the optimum $\rho^*$ such that the derivative provided below can bring the
parameters the rest of the way to the optimum.

\subsection{Main Lemma}
Recall the proposed update direction from Williams (1992) \cite{williams1992}:
\begin{equation}
\Delta\rho \propto \frac{\partial\log P_\rho(a_t|s_t)}{\partial \rho} z_t
\tag{1 revisited}
\end{equation}
Below we will show that the above update direction is too short sighted --- arrives at parameters
$\rho$ that maximizes the one-step increase of the value function --- in other words it asks: ``If I could
only take one move -- what is the best possible move to take?''. The correct equilibrium strategy
would instead say: ``Given that we’re both playing according these two known strategies -- what
is the best possible action to maximize my reward function?''

\begin{lemma}
\begin{enumerate}
\item The original update direction chooses $\rho$ to maximize the value $v$ in one step:
\[
E\!\left[\frac{\partial\log P_\rho(a_t|s)}{\partial \rho}z_t\right]
=E\!\left[\frac{\partial}{\partial \rho}\left(v_\rho(s)-v_\rho(f(s,a_t))\right)\mid s_t=s,\ a_t,\dots,T\sim P_\rho\right].
\]
\item The correct update direction would take into account the full sequence of actions going into
the future:
\[
E\!\left[\frac{\partial v_\rho(s_t)}{\partial \rho}\mid a_t,\dots,T\sim P_\rho\right]
=E\!\left[z_T\sum_{l=t}^T\frac{\partial\log P_\rho(a_l|s_l)}{\partial \rho}\mid a_t,\dots,T\sim P_\rho\right].
\]
\end{enumerate}
\end{lemma}

Proof. We take the expectation of the updating direction referred in \eqref{eq:williams}:
\[
E\!\left[\frac{\partial\log P_\rho(a_t|s)}{\partial \rho}z_t\mid s=s_t,\ a_t,\dots,T\sim P_\rho\right]
=\sum_{a_t\in A(s)}\frac{\partial\log P_\rho(a_t|s)}{\partial \rho}E\!\left[z_{t+1}\mid s_{t+1}=f(s,a_t), a_{t+1},\dots,T\sim P_\rho\right]P_\rho(a_t|s).
\]
\[
=\sum_{a_t\in A(s)}\frac{\partial P_\rho(a_t|s)}{\partial \rho}E\!\left[z_{t+1}\mid s_{t+1}=f(s,a_t), a_{t+1},\dots,T\sim P_\rho\right]
\]
\[
=\sum_{a_t\in A(s)}\frac{\partial P_\rho(a_t|s)}{\partial \rho}v_\rho(f(s,a_t)).
\tag{3}
\]
On examining the value function:
\[
v_\rho(s)=E\!\left[z_t\mid s_t=s,\ a_t,\dots,T\sim P_\rho\right]
=\sum_{a_t\in A(s)}E\!\left[z_{t+1}\mid s_{t+1}=f(s,a_t), a_{t+1},\dots,T\sim P_\rho\right]P_\rho(a_t|s)
\]
\[
=\sum_{a_t\in A(s)}v_\rho(f(s,a_t))P_\rho(f(s,a_t)|s).
\]
We take the derivative of the value function:
\[
\frac{\partial v_\rho(s)}{\partial \rho}
=\sum_{a_t\in A(s)}\left(\frac{\partial v_\rho(f(s,a_t))}{\partial \rho}P_\rho(f(s,a_t)|s)\right)
+\sum_{a_t\in A(s)}v_\rho(f(s,a_t))\left(\frac{\partial P_\rho(f(s,a_t)|s)}{\partial \rho}\right).
\]
If we combine the above result with (3):
\[
E\!\left[\frac{\partial\log P_\rho(a_t|s)}{\partial \rho}z_t\mid s=s_t,\ a_t,\dots,T\sim P_\rho\right]
=\sum_{a_t\in A(s)}\frac{\partial}{\partial \rho}\left(v_\rho(s)-v_\rho(f(s,a_t))\right)P_\rho(f(s,a_t)|s).
\]
Finally, the above result shows that at $s=s_t$,
\[
\frac{\partial\log P_\rho(a_t|s)}{\partial \rho}z_t
\]
is an unbiased estimate of
\[
E\!\left[\frac{\partial}{\partial \rho}\left(v_\rho(s)-v_\rho(f(s,a_t))\right)\mid s_t=s,\ a_t,\dots,T\sim P_\rho\right].
\]
Now we are close to the final result: We can now sum the values
\[
\sum_{l=i}^j\frac{\partial\log P_\rho(a_l|s_l)}{\partial \rho}z_i
\]
going into the future to achieve an unbiased estimate of
\[
E\!\left[\frac{\partial}{\partial \rho}(v_\rho(s_i)-v_\rho(s_j))\mid s_t=s_i,\ a_t,\dots,T\sim P_\rho\right].
\]
We arrive at the correct update direction for the parameters $\rho$:
\[
E\!\left[z_T\sum_{l=t}^T\frac{\partial\log P_\rho(a_l|s_l)}{\partial \rho}\right]
=\frac{\partial}{\partial \rho}\left(v_\rho(s_t)-v_\rho(s_T)\right)
=\frac{\partial v_\rho(s_t)}{\partial \rho}.
\]
The last equality is because $s_T$ is the final state of the game, then $v_\rho(s_T)=r(s_T)$
which does not depend on $\rho$.

\section{Conclusion}
Neural Nets can use reinforcement learning to approximate optimal policy functions that are
intractable to compute and impossible to derive analytically.
Reinforcement learning can be
applied in the vast majority of economic models. We have shown that the conventional update
direction for the parameters in reinforcement learning is too shortsighted. We derived a corrected
update direction that maximizes the reward in any state.
It may be the case that when the policy function $P(a|s)$ is still relatively ``untrained'' and
far from perfect play, looking only one step into the future instead of simulating all the way
towards the end of the game would be better due to the lack of precision with predicting future
play. Here we have shown exactly what the original update direction measures. If the modeler has
qualms about the accuracy of predicting future play, then the modeler may choose to explicitly
discount the corresponding terms in the correct update direction:
\[
\frac{\partial v_\rho(s_t)}{\partial \rho}=E\!\left[z_T\sum_{l=t}^T\frac{\partial\log P_\rho(a_l|s_l)}{\partial \rho}\right].
\]
One should use the full analytic expression we have derived above to check if we have arrived at a
proper local extrema.
Here we have used a general setting to show that the update direction of a reinforcement
learning algorithm should not be as \eqref{eq:williams}. The original update direction
maximized a one step increase in value function.
Instead we show that it is better to look
towards the end of the game to maximize the entire value function itself.

\begin{thebibliography}{12}
\bibitem{bellman1954}
Richard Bellman (1954): \textit{The Theory of Dynamic Programming}.
\bibitem{williams1992}
Ronald J.\ Williams (1992): ``Simple statistical gradient-following algorithms for connectionist
reinforcement learning'' --- \textit{Machine Learning}, Vol.\ 8, 229--256, 1992.
\bibitem{mackay2003}
David J.\ C.\ MacKay (2003): \textit{Information Theory, Inference, and Learning Algorithms}. Cambridge University Press.
\bibitem{kriesel2005}
David Kriesel (2005): \textit{A Brief Introduction to Neural Networks}.
\bibitem{silver2016}
David Silver et al.\ (2016): ``Mastering the game of Go with deep neural networks and tree search'' --- \textit{Nature}.
\bibitem{campbell2002}
Murray Campbell, Joseph Hoane Jr., Feng-hsiung Hsu (2002): \textit{Deep Blue} --- \textit{Artificial Intelligence}. Vol.\ 134, Nos.\ 1--2, 57--83.
\bibitem{schaeffer1992}
Jonathan Schaeffer et al.\ (1992): ``A World Championship Caliber Checkers Program'' --- \textit{Artificial Intelligence}. Vol.\ 53, 273--289.
\bibitem{buro1999}
Michael Buro (1999): ``From simple features to sophisticated evaluation functions'' --- 1st
International Conference on Computers and Games.
\bibitem{tesauro1996}
Gerald Tesauro and Gregory (1996): ``On-line policy improvement using Monte-Carlo tree
search'' --- Advances in Neural Information Processing, 1068--1074.
\bibitem{sheppard2002}
Brian Sheppard (2002): ``World Championship Caliber Scrabble'' --- \textit{Artificial Intelligence}. 134, 241--275.
\bibitem{spence1973}
Michael Spence (1973): ``Job Market Signalling'' --- \textit{Quarterly Journal of Economics}.
\bibitem{hornik1989}
Kurt Hornik, Maxwell Stinchcombe, Halber White (1989): ``Multilayer Feedforward Networks are Universal Approximators''.
\end{thebibliography}

\end{document}
